\chapterrule
\section{反向传播究竟在传什么？}
\phantomsection\label{concept:backprop}
\label{intro:backpropagation}

多层网络的参数通过后续计算逐步影响最终损失。
若为每个参数都重新展开整条计算链，会做许多重复运算。
\textbf{反向传播}要解决的，是复用同一张计算图里的局部导数，
一次高效求出所有参数的梯度。

\subsection{从一条计算链理解链式法则}

先看一台三级机器：
\[
x\xrightarrow{\times\,w_1}h
\xrightarrow{\times\,w_2}\widehat{y}
\xrightarrow{\text{算损失}}L.
\]
令 $h=w_1x$、$\widehat{y}=w_2h$、$L=\tfrac12(\widehat{y}-y)^2$。
这里只有一个样本，所以不再像前几章那样对三点求平均；
系数 $\tfrac12$ 只为让导数 $\partial L/\partial\widehat{y}=\widehat{y}-y$ 更简洁。
现在要问：调整 $w_1$ 一点点，最后的损失 $L$ 会变多少？

最直接的做法是把整个函数展开，但神经网络可以有数十乃至更多层。链式法则沿着计算路径，把各步的\textbf{局部变化率}相乘：
\begin{equation}
\boxed{\frac{\partial L}{\partial w_1}
=\frac{\partial L}{\partial\widehat{y}}
\frac{\partial\widehat{y}}{\partial h}
\frac{\partial h}{\partial w_1}.}
\label{eq:chain-three}
\end{equation}
这叫\textbf{链式法则}（chain rule）。

\subsection{亲手算一次两层网络的梯度}
\label{sec:backprop-weight-gradients}

取 $x=2$、真实目标 $y=5$、初始权重 $w_1=1$、$w_2=2$。
前向过程：
\[
h=1\times2=2,\quad
\widehat{y}=2\times2=4,\quad
L=\tfrac12(4-5)^2=\tfrac12.
\]

从输出向输入反向算：
\[
\frac{\partial L}{\partial\widehat{y}}
=\widehat{y}-y=-1,\qquad
\frac{\partial\widehat{y}}{\partial h}=w_2=2,\qquad
\frac{\partial h}{\partial w_1}=x=2.
\]
所以
\[
\boxed{\frac{\partial L}{\partial w_1}
=(-1)\times2\times2=-4.}
\]
而对于最后一层权重：
\[
\frac{\partial L}{\partial w_2}
=\frac{\partial L}{\partial\widehat{y}}
\frac{\partial\widehat{y}}{\partial w_2}
=(-1)\times h=-2.
\]
两条梯度都为负：略增 $w_1$ 或 $w_2$，预测 $4$ 都会向目标 $5$ 靠近。
\hyperref[sec:end-to-end-step]{本章最后一节}给这个网络加上两项偏置，
只需补算偏置的梯度，再让四个参数一起更新。

% The Figure structure intentionally includes the non-floating caption.
\tagstructbegin{tag=Figure,alttext={两层神经网络的前向计算和参数梯度回传路线}}
\begin{center}
\begin{tikzpicture}[>=Latex,
 block/.style={draw=linegray,fill=paper,rounded corners=2pt,minimum width=33mm,
 minimum height=14mm,align=center,font=\small},
 param/.style={draw=flowcolor,fill=linegray!25,rounded corners=2pt,
 minimum width=20mm,minimum height=8mm,align=center,font=\small}]
\node[block] (h) at (0,0) {$h=w_1x$\\前向：$1\times2=2$};
\node[block] (out) at (4.3,0) {$\widehat{y}=w_2h$\\前向：$2\times2=4$};
\node[block] (loss) at (8.6,0) {$L=\frac12(\widehat{y}-5)^2$\\前向：$L=0.5$};
\node[param] (w1) at (0,1.6) {$w_1=1$};
\node[param] (w2) at (4.3,1.6) {$w_2=2$};
\draw[->,flowcolor,thick] (w1)--(h);
\draw[->,flowcolor,thick] (w2)--(out);
\draw[->,flowcolor,thick] (h)--node[above,font=\scriptsize]{$h=2$}(out);
\draw[->,flowcolor,thick] (out)--node[above,font=\scriptsize]{$\widehat{y}=4$}(loss);
\node[font=\footnotesize,flowcolor,align=center] (g1) at (0,-1.65)
 {$\frac{\partial L}{\partial h}=-2$\\$\frac{\partial L}{\partial w_1}=-4$};
\node[font=\footnotesize,flowcolor,align=center] (g2) at (4.3,-1.65)
 {$\frac{\partial L}{\partial\widehat{y}}=-1$\\$\frac{\partial L}{\partial w_2}=-2$};
\node[font=\footnotesize,flowcolor,align=center] (g3) at (8.6,-1.65)
 {从损失开始反传\\$\frac{\partial L}{\partial\widehat{y}}=4-5$};
\draw[->,flowcolor,thick] (g3.west)--(g2.east);
\draw[->,flowcolor,thick] (g2.west)--(g1.east);
\node[font=\footnotesize,align=center] at (4.3,-2.65)
 {以 $\eta=0.05$ 更新：$w_1:1\to1.2$，$w_2:2\to2.1$（偏置更新见正文）};
\end{tikzpicture}
\end{center}
\begin{center}
\captionof{figure}{上方是前向值和参数节点，下方是反向梯度及一次参数更新}
\label{fig:backprop}
\end{center}
\tagstructend


\begin{idea}
反向传播（backpropagation）传的不是文字、答案或原始输入，而是损失对中间变量的梯度。
前向时从 $x$ 算到 $L$；求梯度时从 $\partial L/\partial\widehat{y}$ 出发，沿同一条链倒着走，这就是“反向”的含义。
后面的局部变化率先算出来，前面的参数直接复用，避免每求一个权重都从头展开整条函数。
\end{idea}\phantomsection\label{route:backprop-end}

\subsection{如果中间有 ReLU，怎样求导？}

先只看 ReLU 这一步：输入记作 $z$，输出记作 $h$，即
\[
h=\relu(z)=\max(0,z).
\]
求导就是问：\textbf{输入稍微变化时，输出跟着变化多少？}
图~\ref{fig:relu-function} 中的两段直线，可以分开来看。

\textbf{输入为正时，$h=z$。}
例如 $z$ 从 $2$ 增加到 $2.1$，$h$ 也从 $2$ 增加到 $2.1$。
输出变化量除以输入变化量是 $0.1/0.1=1$，所以这一段的导数为 $1$。

\textbf{输入为负时，$h=0$。}
例如 $z$ 从 $-2$ 增加到 $-1.9$，$h$ 始终为 $0$。
输出变化量除以输入变化量是 $0/0.1=0$，所以这一段的导数为 $0$。
只要输入的变化留在同一段内，把变化量缩得再小，这两个比例也不变。因此
\[
\frac{\partial h}{\partial z}=\relu'(z)=
\begin{cases}
0,&z<0,\\
1,&z>0.
\end{cases}
\]

\textbf{恰好在 $z=0$ 时，两段在这里形成折角。}
从左侧靠近，斜率是 $0$；从右侧靠近，斜率是 $1$。
两侧不一致，所以这一点没有通常意义下的导数。
但反向传播程序仍需要一个数来完成下面的乘法。
本书把零点按负数一侧处理，约定：\textbf{若输入恰好为 $0$，反传经过 ReLU 时也乘 $0$。}
这是折角处采用的计算规则，并不是求出了该点的导数。

现在把 ReLU 接回网络：
$z=w_1x+b_1$、$h=\relu(z)$、$\widehat{y}=w_2h+b_2$。
先从输出端算出 $\partial L/\partial h$；反传经过 ReLU 时，乘上刚求出的局部导数：
\[
\frac{\partial L}{\partial z}
=\frac{\partial L}{\partial h}\frac{\partial h}{\partial z},\qquad
\frac{\partial L}{\partial w_1}
=\frac{\partial L}{\partial z}\frac{\partial z}{\partial w_1}
=\frac{\partial L}{\partial z}\,x.
\]

沿用前面的 $x=2,y=5,w_1=1,w_2=2$，两项偏置取 $0$。
此时 $z=2,h=2,\widehat{y}=4$，所以
\[
\frac{\partial L}{\partial h}=(4-5)\times2=-2,\qquad
\frac{\partial L}{\partial w_1}=(-2)\times\underbrace{1}_{\text{ReLU 导数}}\times2=-4.
\]
只把输入改成 $x=-2$，目标和参数不变，则 $z=-2,h=0,\widehat{y}=0$，此时
\[
\frac{\partial L}{\partial h}=(0-5)\times2=-10,\qquad
\frac{\partial L}{\partial w_1}=(-10)\times\underbrace{0}_{\text{ReLU 导数}}\times(-2)=0.
\]
虽然预测仍有误差，乘积中的 ReLU 导数为 $0$，使这个样本经这条路径传给 $w_1$ 的梯度为 $0$。
这就是“梯度被截断”的含义：负数段的输出保持不变，前面参数的微小变化暂时无法沿这条路径影响损失。
其他不经过这个 ReLU 的路径仍可能提供梯度。

\subsection{反向传播负责求梯度，优化器负责更新}

\begin{center}
\tagpdfsetup{table/header-rows={1}}
\begin{tabularx}{\linewidth}{@{}p{29mm}X@{}}\toprule
概念 & 它负责什么\\\midrule
前向传播 & 给定参数，算输出和损失\\
反向传播 & 高效计算损失对每个参数的梯度\\
梯度下降/优化器 & 根据梯度修改参数\\
\bottomrule\end{tabularx}
\end{center}

更复杂的计算图会有分叉和相加：同一个参数影响损失的多条路径需要把各条路径的梯度贡献相加。自动微分框架正是在计算图上应用这个规则。

\begin{warningbox}
每个训练步都用同一组旧参数完成前向与反向传播，算出全部梯度，再由优化器统一更新所有参数。
不要反传到某一层就立刻修改那一层，否则同一步会混用新旧权重。
\end{warningbox}

\subsection{把前面的数学连成一次真正的学习}
\label{sec:end-to-end-step}

前面已经会用函数前向计算、用损失衡量误差，并沿两层计算链反传梯度。
把这些步骤连起来，就是一次参数学习。
继续沿用\hyperref[sec:backprop-weight-gradients]{本章前面}的 $x=2,y=5,w_1=1,w_2=2$，这次添加两项偏置，
从\textbf{预测}算到\textbf{四个参数同时更新后的再预测}。

\subsubsection{先写一个只有一个隐藏单元的小网络}

设
\begin{equation}
z=w_1x+b_1,\quad h=\relu(z),\quad
\widehat{y}=w_2h+b_2,\quad
L=\frac12(\widehat{y}-y)^2.
\label{eq:small-network-full}
\end{equation}
从左到右分别是输入、第一层线性变换、激活函数、第二层线性变换、损失。虽然只有一个隐藏单元，但计算方式与更大 MLP 的局部结构相同。
初始参数取 $w_1=1,b_1=0,w_2=2,b_2=0$。

\subsubsection{前向与反向：只需补算两项偏置梯度}

两项偏置都为零，前向结果与前面相同：
\[
z=1\times2+0=2,\quad h=\relu(2)=2,\quad
\widehat{y}=2\times2+0=4,\quad L=\tfrac12.
\]
反向时仍有 $\partial L/\partial\widehat{y}=\widehat{y}-y=-1$。
因为 $z=2>0$，$\relu'(2)=1$，ReLU 这一段原样传递梯度，
两项权重梯度仍是 $\partial L/\partial w_1=-4$ 与 $\partial L/\partial w_2=-2$。
偏置对自己所在那一层的局部导数都是 $1$，所以
\[
\frac{\partial L}{\partial b_2}=(-1)\times1=-1,\qquad
\frac{\partial L}{\partial b_1}=(-1)\times w_2\times\relu'(2)\times1=(-1)\times2\times1=-2.
\]
若 $z$ 为负，$\relu'(z)=0$，$w_1$ 与 $b_1$ 的梯度就会同时变成零。

\subsubsection{四个参数一起更新，再做一次预测}

使用学习率 $\eta=0.05$：
\[\begin{aligned}
w_1'&=1-0.05(-4)=1.2,&
b_1'&=0-0.05(-2)=0.1,\\
w_2'&=2-0.05(-2)=2.1,&
b_2'&=0-0.05(-1)=0.05.
\end{aligned}\]

用新参数重新预测：
\[\begin{aligned}
z'&=1.2\times2+0.1=2.5,\\
h'&=\relu(2.5)=2.5,\\
\widehat{y}'&=2.1\times2.5+0.05=5.30,\\
L'&=\frac12(5.30-5)^2=0.045.
\end{aligned}\]
损失从 $0.5$ 降到 $0.045$，这次更新改善了该样本的预测。

\subsubsection{下一步：如果有许多输入数字怎么办？}

这一章仍只有单个 $x$ 与 $h$。若同时处理面积、房龄、位置等特征，
或有许多隐藏单元，我们仍然使用同一个学习循环，
但逐一写所有乘法和导数会非常繁琐。
后面会用\textbf{向量和矩阵}表达同样的计算，而不改变学习机制本身。

% The earlier backpropagation lesson belongs to this same numbered chapter.
% Keep both complete exercise sets after section 5.5, not before it.
\chapterchecktasks{backpropagation}{基础与贯通复算}{
\item \textbf{计算：}设 $h=w_1x,\widehat{y}=w_2h$。当 $x=3,w_2=2$ 时，按链式法则求 $\partial\widehat{y}/\partial w_1$。\chapterinlinehint{沿输出到隐藏层、隐藏层到权重逐段求导。}
\item \textbf{预测：}在 $h=\relu(z)$ 且 $z=-2$ 时，从输出侧传来的梯度经过这条 ReLU 分支传向 $z$ 后还剩多少？\chapterinlinehint{先看该输入处 ReLU 的局部导数。}
\item \textbf{找错：}同学把 $\partial L/\partial h$ 直接当作 $\partial L/\partial w_1$。缺了哪一个局部导数？\chapterinlinehint{补上隐藏变量对权重的局部导数。}

\item \textbf{流程：}把一个训练步分成“前向／损失／反传／更新”四步；哪一步读取训练目标，哪一步改变参数？\chapterinlinehint{先标出损失与目标的连接，再找赋值更新的位置。}
\item \textbf{手算：}$x=2,y=5,w_1=1,w_2=2,b_1=b_2=0$，$h=\relu(w_1x+b_1)$。求 $\widehat{y}$ 与 $\tfrac12(\widehat{y}-y)^2$。\chapterinlinehint{先做隐藏层前向，再算预测与目标的平方差。}
\item \textbf{手算：}沿用上题参数，只把输入改成 $x=-1$。求 $z,\widehat{y}$ 与 $\partial L/\partial\widehat{y}$，再指出四个参数中哪些梯度为零。\chapterinlinehint{先判断 $z$ 的正负，再看 ReLU 能否传递梯度。}
}


\par\smallskip
训练的基本循环是：预测 $\to$ 算损失 $\to$ 求梯度 $\to$ 更新参数。
这个例子说明了参数怎样被更新；至于新样本是否预测得准，还涉及泛化问题。
下一章解释什么叫泛化，以及怎样识别模型把样本噪声当成规律。\phantomsection\label{route:backprop-all-end}
